\section{Conclusion}
\label{sec:conclusion}

We find that quantization induced instruction following failures are highly
localized, arising from disruption in a small subset of attention heads.
To address this, we introduce Deep Attention Stimulation (DAS), a training free
intervention that restores instruction following by compensating degraded
attention head outputs.
DAS recovers language and formatting constraints under moderate stimulation,
while excessive stimulation destabilizes generation.
These results point to a new direction for targeted and interpretable
post quantization model repair.

% We introduce \emph{Diagnostic Attention Stimulation (DAS)}, a training-free
% approach for recovering instruction-following behavior in quantized large
% language models through targeted attention head interventions.
% Our analysis shows that quantization induces localized, sparse disruptions in a
% small subset of critical attention heads rather than uniform degradation.
% By selectively compensating these disrupted heads with moderate stimulation
% ($\alpha=5$), DAS recovers instruction-following behavior in most qualitative
% failure cases, while stronger stimulation leads to degeneration.
% These findings suggest that quantization-induced failures are diagnosable and
% correctable through targeted, interpretable interventions, pointing to a new
% direction for post-quantization model repair.